\chapter{Discussion, Limitations, and Recommendations}
\section{Introduction}
The main objective of this chapter is to summarize an overview of this research, evaluating the limitations of the project as well as offering recommendations and suggestions on areas that need improvement and further research. As discussed in the previous chapters, the main objectives of this project were to build a Bayesian Framework model for quantifying causal relationships of operational risk. Operational risk posed a significant challenge in the banking sector, given the rapid transformation in digital technology across different financial institutions. This thesis provided clear and understandable recommendations based on the output and results from the analysis. This has particularly provided benefits to the researcher, the University itself, and the main target audience, which includes the banks, Insurance providers, and other financial institutions. The main goal is to reduce and curb the costs associated with operational risk. 

\section{Discussion} 
The Bayesian Inference analysis was done using a simulated dataset spanning from 1980 to 2024, and the results of the model produced were accurate and fit. The model was tested and validated using a real industry dataset for operational risk. The variables in question were seven, which indicate the main key risk indicators according to the Basel II Standard for banking supervision. Causal discovery and the relationship between the variables were examined as well. The following results and discoveries were made, alongside the recommendations.
\subsection{Objective 1, Model Construction}
From the findings, the model established the hierarchical structure of the key risk indicators, with the most prominent risk factors at the top, the primary risk sources (Internal Fraud, External Fraud, ICT Failure), risk categories (Employment Practices, Execution and Process Management, Clients/Products and Business Practices, Disaster), and the central operational risk node. Other risk factors follow, and the main parent node at the centre is the Total Operational Risk Loss. The Directed Acyclic Graph was constructed to show the dependencies of the risk factors. The findings reveal that there is a direct link between nodes, which indicates that the nodes directly affect each other. 

However, other nodes do not have a direct link with others but contribute directly to the total loss. ICT and System failure create loopholes to internal fraud within the organisation. Natural disasters cause external fraud, phishing, and cyber threats, just to mention a few. External Fraud gives rise to execution delivery, Client's products, and Business practices, and lastly to total loss. From the model analysis, we have found that Clients, Products, and Business practices have the highest percentage of the total impact on the total operational risk. These are the notable results obtained from the model, suggesting a robust Bayesian Network for risk Quantification, which Banks, Insurance, and other financial institutions can use as a risk management tool for efficient risk management.

\subsection{Objective 2: Network Quantification}
Objective two: focused on quantifying the Bayesian Network model of operational risk. The Prior distributions were assigned to each risk factor based on expected behavior, and inference was performed using the Bayesian maximum likelihood estimation method. The Posterior distributions showed high probability estimates (0.91–0.98) with varying certainty. External fraud was estimated precisely, while Employment Practices showed greater uncertainty. Density plots revealed diverse beta distributions, and convergence diagnostics (trace plots and R-hat = 1) confirmed reliable sampling. A key challenge was handling uncertainty in factors with flatter or more dispersed distributions, such as internal fraud.

\subsection{Objective 3: Model Validation}
The validation process, including scenario analysis testing and cross-validation, established that the model's structure accurately represents causal relationships, with posterior distributions predominantly driven by scenario-specific data patterns rather than modeling assumptions. This comprehensive approach provides a reliable framework for operational risk assessment that balances prior information with observed historical data, supporting evidence-based decision-making.

\section{Conclusion} 
This research established a framework for quantifying, validating, and interpreting the causal relationships between operational risk factors in organizations, particularly in the field of insurance and banking sectors. The main objectives were successfully achieved through a structured approach. The model identified and quantified the impact of key risk attributes for the identified key risk indicators. By estimating the probability of various risk scenarios and validating the framework through sensitivity and scenario analysis, this dissertation demonstrated a strong Bayesian Model that can serve as a valuable tool for risk management.

A comprehensive BBN diagram was developed, providing a hierarchical view of risk factors. The model visually illustrates how risks propagate through an organization and how different risk factors are interdependent. Key findings from this framework include the identification of primary risk sources (e.g., Internal Fraud and ICT Failure) and the ratio of their impact on total Loss. The model’s probabilistic approach allows decision-makers to assess and prioritize risks more effectively. Statistical validation confirmed the robustness of these estimates, with significant p-values $(p < 0.001)$) indicating strong confidence in the network’s ability to represent causal relationships.

The BBN Model was validated through scenario testing and sensitivity analysis. The results revealed consistent patterns across different risk scenarios, with minimal variation in key risk parameters, confirming the model's robustness. Sensitivity analysis further ensured that the BBN could maintain its predictive power even under varying prior assumptions, suggesting that the model is reliable for decision-making under different risk conditions.

Through the Directed Acyclic Graph (DAG) analysis, this research uncovered direct and indirect causal links between risk factors. For instance, system failure creates opportunities for internal fraud, while natural disasters trigger external fraud and other risks. These findings emphasize the interconnected nature of operational risks and highlight the importance of a comprehensive approach to risk management.

The study's findings underscore the strategic value of using Bayesian Networks to guide risk management practices. By identifying the most significant risk factors and their relationships, organizations can prioritize risk mitigation efforts in high-impact areas such as client onboarding processes, product governance, and fraud detection systems. This framework offers a proactive approach to managing risks.

This research demonstrates that BBN offers a powerful tool for operational risk modeling, providing a structured framework to understand, quantify, and manage complex risk relationships. By integrating statistical analysis, expert judgment, and scenario testing, the Bayesian Network model introduces a more sophisticated, dynamic, and data-driven approach to managing risk. This approach aligns with contemporary regulatory expectations, such as Basel III and ICAAP, and provides organizations with the tools they need to navigate the evolving landscape of operational risk.
\section{Recommendations}
Organizations, should consider adopting Bayesian Network models for operational risk management. These models offer a more comprehensive understanding of risk factors, enabling better decision-making and prioritization of risk mitigation strategies.
Industry practitioners should incorporate scenario testing and sensitivity analysis into their risk management practices. These techniques can help validate risk models and ensure robustness under different operational conditions.

Since operational risk dynamics evolve, organizations should invest in continuous data collection efforts, particularly in high-risk areas, to refine their risk models and minimize uncertainties in predictions.

Additional variables, such as market fluctuations and geopolitical risks, could be incorporated into future models to expand the scope of operational risk modeling. This would help create a more robust and adaptable model for different sectors.

The university should continue to support advanced risk management methodologies, including Bayesian Neural Networks, by providing resources and opportunities for interdisciplinary collaboration with the insurance, banking, and financial sectors. Encouraging collaboration with industry partners can provide valuable insights and practical applications for academic research, allowing students and researchers to test and refine theoretical models in real-world settings.
\section{Further Research and Study}

Future research should focus on validating the Bayesian Network framework using extensive real-world datasets from different financial institutions. This would improve the external validity and robustness of the model across diverse operational environments. A valuable extension of this study recommends the integration of Dynamic Bayesian Models, which capture the evolution of operational risk over time.  This would allow modelling of temporary risk factors that might occur more frequently than expected over some time. This offers a flexible modelling approach to forecast the future cost of operational risks. 

As an Actuarial science student, I would recommend the integration of operational risk models with credit and market risk models to enhance strategic decision-making. Future work may also consider applying more advanced causal discovery algorithms, such as deep learning, Greedy Equivalence Search (GES), Fast Causal Inference (FCI), or constraint-based methods to improve the accuracy and interpretability of the causal relationships within the network. Due to more sophisticated algorithms, future research should explore Random Forest and Neural Networks to improve the predictive performance. 

Bayesian Network analysis is not only a tool for quantifying risk in financial institutions alone, we can diversify to non-financial sectors like healthcare, manufacturing, or energy, where operational risk is very high given the complex nature of their systems. One major concern in this research is that the Data for risk analysis is sensitive and very scarce. Therefore, for future research, I would suggest incorporating expert judgment and elicitation, like the Delphi method, to enhance the accuracy and fit of the model.


